\documentclass[11pt,a4paper]{article}
\usepackage[utf8]{inputenc}
\usepackage[T1]{fontenc}
\usepackage{amsmath,amssymb,amsthm}
\usepackage[margin=2.5cm]{geometry}
\usepackage{xcolor}
\usepackage{listings}
\usepackage{booktabs}
\usepackage[colorlinks=true,linkcolor=blue!60!black,citecolor=blue!60!black,urlcolor=blue!60!black]{hyperref}

\definecolor{leanblue}{RGB}{43,101,236}
\definecolor{leangreen}{RGB}{11,102,35}
\definecolor{warnred}{RGB}{170,30,30}

\lstdefinelanguage{lean}{
  morekeywords={def, theorem, lemma, by, trivial, exact, intro, apply,
    noncomputable, namespace, end, Prop, Real, Complex, open, import, sorry,
    axiom, structure, where},
  sensitive=true,
  morecomment=[l]{--},
  morecomment=[s]{/-}{-/},
  morestring=[b]",
}
\lstset{
  basicstyle=\ttfamily\small,
  frame=single,
  columns=fullflexible,
  keepspaces=true,
  commentstyle=\color{leangreen},
  keywordstyle=\color{leanblue}\bfseries,
  language=lean,
  numbers=left,
  numberstyle=\tiny\color{gray}
}

\theoremstyle{plain}
\newtheorem{conjecture}{Conjecture}
\newtheorem{proposition}{Proposition}
\theoremstyle{definition}
\newtheorem{definition}{Definition}
\theoremstyle{remark}
\newtheorem{remark}{Remark}

\newcommand{\tierA}{\textbf{Tier A}}
\newcommand{\tierB}{\textbf{Tier B}}
\newcommand{\tierC}{\textbf{Tier C}}

\title{\vspace{-1.5cm}\textbf{Ramanujan Neuro-Symbolic Mathematics:\\
A Research Program for Dual-Scale Geometry,\\
\texorpdfstring{$K3\times T^{2}$}{K3 x T2} Compactifications, and the Structure of Singular Cascades}\\[0.4em]
\large Version 2 --- position paper; contains no new proven theorems}
\author{\textbf{Xavier Callens}\\ \small SocrateAI Lab}
\date{August 9, 2026}

\begin{document}
\maketitle

\begin{center}
\fbox{\parbox{0.92\textwidth}{\small
\textbf{Epistemic status.} This is a \emph{program proposal}, not a results
paper. Every claim below carries a tier label defined in
Section~\ref{sec:protocol}: \tierA{} (machine-checked), \tierB{} (established
literature, cited), \tierC{} (conjecture or heuristic). At the time of
writing, the program's central dictionary contains \emph{no} Tier~A entries.
Version~2 supersedes and corrects a previous draft whose Lean~4 listings were
vacuous; see the Errata in Section~\ref{sec:errata}.}}
\end{center}

\begin{abstract}
\noindent We outline a research program that seeks \emph{exact algebraic
surrogates} --- drawn from $q$-series, mock modular forms, and the arithmetic
of $K3$ surfaces --- for the analytic estimates that govern regularized fluid
dynamics and related dual-scale geometries. The program couples (i) a
heuristic search engine (\textsc{Rama}) over Eulerian $q$-series, (ii) shadow
completion in the sense of Zwegers, and (iii) formalization in the Lean~4
proof assistant under a strict audit discipline. We state the program's five
central conjectures precisely, grade every claim within a three-tier
epistemic system, and attach to each conjecture an explicit falsification
protocol and a formal target. This version withdraws the ``machine-verified''
claims of the previous draft, whose theorem statements were of type
\texttt{True} and therefore certified nothing. The contribution of the
present document is the honest formulation of the conjectures, the audit
rules under which any future claim of proof must be certified, and a
milestone plan whose outcomes are decidable.
\end{abstract}

\section{The Proposal}\label{sec:proposal}

For most of its history, rigorous mathematics has proceeded bottom-up:
deriving $C$ from $A$ and $B$. Ramanujan's working style --- documented most
starkly in the \emph{Lost Notebook}~\cite{Narosa1988,AndrewsBerndt} --- was
different: identities were recorded as endpoints, with the connective proofs
supplied by later generations. We take that style as \emph{design
inspiration}, not as evidence for any particular cognitive model (see
Section~\ref{sec:lostnotebook}), and propose a three-stage pipeline:

\begin{enumerate}
\item \textbf{The heuristic engine (\textsc{Rama}).} Mathematical candidate
generation is treated as local search over a symbolic space of Eulerian
$q$-series, guided by an energy functional
$E=\alpha C+\beta I+\gamma D$ (complexity, inconsistency, distance-to-anchor
terms). \tierC{}: the functional is a design choice; no optimality claim is
made.
\item \textbf{The shadow bridge.} Candidate mock-symmetries are completed to
harmonic Maass forms by computing their non-holomorphic period integrals,
following Zwegers~\cite{Zwegers2002} (e.g.\ the weight-$3/2$ shadow
$\eta(\tau)^{3}$ arising in the $K3$ elliptic-genus decomposition). \tierB{}
as a technique; \tierC{} for any specific new identity produced by the
engine until certified.
\item \textbf{The formal lock.} Surviving statements are formalized in
Lean~4 under the audit rules of Section~\ref{sec:protocol}.
\end{enumerate}

A caveat governs the entire pipeline. \emph{A proof kernel verifies exactly
what is stated, and nothing more.} A theorem whose statement is \texttt{True}
is verified and worthless. The value of stage (iii) is therefore bounded by
the strength of the statements we manage to formalize, not by the presence
of a green checkmark. The previous draft violated this principle; the
present one is organized around it.

\section{Epistemic Protocol}\label{sec:protocol}

Every assertion in this program carries exactly one grade.

\begin{center}
\begin{tabular}{@{}lll@{}}
\toprule
Tier & Meaning & Admission criterion \\
\midrule
A & Machine-checked & Lean~4 proof; \texttt{\#print axioms} clean; no \texttt{sorry} \\
B & Established & Peer-reviewed literature, cited inline \\
C & Conjecture / heuristic & Everything else, however plausible \\
\bottomrule
\end{tabular}
\end{center}

\noindent The following rules are binding on all program artifacts.

\begin{description}
\item[R1.] No \texttt{axiom} declarations. (The current
\texttt{DualScale.lean} baseline satisfies this: zero axioms, thirteen audit
certificates.)
\item[R2.] \texttt{sorry} marks a target as \textsc{open}. It is never
counted as verified, and its presence must be reported.
\item[R3.] Theorem statements whose type is \texttt{True}, or is otherwise
trivially inhabited independently of the mathematical content, are
prohibited.
\item[R4.] Every Tier~B claim cites its source. Uncited claims default to
Tier~C.
\item[R5.] Numerical claims require exact-arithmetic \textsc{pass/fail}
certificates over $\mathbb{Q}$, with inputs locked to their literature
sources --- the discipline already used by the program's moduli-map checker,
which produced \textsc{pass}(40) for the (K-$s_{7}$, E-$s_{7/2}$) pairing and
\textsc{fail} for a cross-level pairing.
\end{description}

\section{The Conjectural Dictionary}\label{sec:dictionary}

The program's core object is a proposed dictionary between the
$\sqrt{\alpha'}$-regularized dynamics of an incompressible flow on a
macroscopic manifold and arithmetic structures carried by a quantum fiber.
Each entry below is presented in four parts: the established anchor
(\tierB{}), the conjecture (\tierC{}), the formal target, and the
falsification test.

\subsection{Hydrodynamic limits and Ramanujan's sums of tails}
\label{sec:ns}

\paragraph{Anchor (\tierB{}).} Ramanujan's ``sums of tails'' identities give
exact closed forms for the defect between certain $q$-series and their
partial sums; a modern treatment, including connections to values of
$L$-functions, is Andrews--Jim\'enez-Urroz--Ono~\cite{AJO2001}. On the
analytic side, each Fourier--Galerkin truncation of the 3D Navier--Stokes
system at frequencies $|k|\le(\alpha')^{-1/2}$ is a finite-dimensional ODE
system and is globally regular by classical theory; blow-up criteria for the
untruncated system are governed by results of Beale--Kato--Majda
type~\cite{BKM1984}.

\paragraph{Conjecture (\tierC{}).}
\begin{conjecture}[Hypothesis U]\label{conj:U}
There exists $C>0$, independent of $\alpha'$, such that the enstrophy
density of the $\sqrt{\alpha'}$-truncated flow satisfies
$\sup_{t}\;\alpha'\!\cdot\!\mathcal{E}_{\alpha'}(t)\le C$ for all
sufficiently small $\alpha'>0$, and the truncation defect of the energy flux
admits an exact sums-of-tails representation.
\end{conjecture}

\begin{remark}
The substantive content is \emph{uniformity in $\alpha'$}. Regularity of
each truncated system is classical and proves nothing about Navier--Stokes;
and even granting Conjecture~\ref{conj:U}, the passage to the
$\alpha'\to 0$ limit requires a compactness argument that is itself open.
No claim of having ``solved'' or ``bypassed'' the Navier--Stokes problem is
made here, and the corresponding claim in the previous draft is withdrawn.
\end{remark}

\paragraph{Falsification.} Exhibit, in DNS surrogates, growth of
$\alpha'\cdot\mathcal{E}_{\alpha'}$ along a sequence $\alpha'\to 0$; or
prove that the flux-defect series is not of sums-of-tails type.

\subsection{Spectral gaps and Ramanujan graphs}\label{sec:gap}

\paragraph{Anchor (\tierB{}).} Deligne's proof of the Ramanujan--Petersson
conjecture~\cite{Deligne1974} gives $|\tau(p)|\le 2p^{11/2}$ for the
Ramanujan $\tau$-function. Lubotzky--Phillips--Sarnak~\cite{LPS1988} used
bounds of this type to construct $k$-regular Ramanujan graphs, whose
non-trivial adjacency eigenvalues satisfy $|\lambda|\le 2\sqrt{k-1}$ --- the
optimal spectral gap by Alon--Boppana.

\paragraph{Conjecture (\tierC{}).}
\begin{conjecture}\label{conj:gap}
The triad-interaction graph of sub-cutoff modes, with edge weights induced
by the fiber arithmetic, is asymptotically Ramanujan. Consequently,
resonant nonlinear echoes decay at the optimal mixing rate set by the
spectral gap.
\end{conjecture}

\begin{remark}
Deligne's bound is a theorem about modular forms; that it \emph{governs} the
interaction graph of a truncated flow is precisely what is conjectured, not
established. Nothing is ``erased'' or ``weaponized'' by citation alone.
\end{remark}

\paragraph{Falsification.} Compute the adjacency spectra of the induced
graphs for small primes $p\in\{2,3,5\}$ under Rule~R5. A single certified
eigenvalue outside the Alon--Boppana window refutes the conjecture in that
regime.

\subsection{Vortex-direction constraints and elliptic integrals}
\label{sec:cfm}

\paragraph{Anchor (\tierB{}).} Constantin--Fefferman~\cite{CF1993} showed
for 3D Navier--Stokes that Lipschitz-type control of the vorticity direction
$\xi=\omega/|\omega|$ in regions of high vorticity precludes blow-up;
Constantin--Fefferman--Majda~\cite{CFM1996} developed the analogous
geometric constraints for 3D Euler.

\paragraph{Conjecture (\tierC{}).}
\begin{conjecture}\label{conj:cfm}
For the truncated flow, the Lipschitz modulus $\|\nabla\xi\|_{\infty}$
equals, up to an explicit constant, the truncation angle
$\alpha<\pi/2$ of an incomplete elliptic integral attached to the fiber,
with the angle determined by the T-dual cutoff.
\end{conjecture}

\begin{remark}
At present there exists \emph{no derivation} linking the modular
parametrization of incomplete elliptic integrals to the direction field of a
Navier--Stokes vortex. The correspondence is proposed on structural grounds
(both objects arise as truncations of a scale symmetry) and stands or falls
with Milestone~M3 of Section~\ref{sec:milestones}.
\end{remark}

\paragraph{Falsification.} Derive the leading-order relation and compare, in
exact arithmetic where possible, against the CF criterion constants; a
mismatch at leading order refutes the identification.

\subsection{Criticality, continued fractions, and the singular set}
\label{sec:phase}

\paragraph{Anchor (\tierB{}).}
Caffarelli--Kohn--Nirenberg~\cite{CKN1982} proved that for suitable weak
solutions of 3D Navier--Stokes the singular set has vanishing
one-dimensional parabolic Hausdorff measure; in particular its dimension is
at most~1. Separately, the Rogers--Ramanujan continued fraction exhibits
genuinely discrete limiting behavior on the unit circle: at certain roots of
unity it diverges while possessing finitely many subsequential limits
\cite{BowmanMcLaughlin2004,AndrewsBerndt}.

\paragraph{Conjecture (\tierC{}).}
\begin{conjecture}\label{conj:dim0}
(i) The singular set of suitable weak solutions has Hausdorff dimension
zero. (ii) As $\alpha'\to 0$, the truncated flow selects among finitely many
non-communicating topological states indexed by continued-fraction limit
classes, rather than diverging to a single blow-up profile.
\end{conjecture}

\begin{remark}
Part~(i) is strictly \emph{stronger} than the known CKN bound and is open;
the previous draft presented it as ``formally proved,'' which was false.
Part~(ii) transfers a theorem about special functions to a statement about
PDE dynamics; the transfer is the conjecture.
\end{remark}

\paragraph{Falsification.} Any construction of a singular set of positive
dimension refutes (i); numerical continuation exhibiting a continuum of
limit states as $\alpha'\to 0$ refutes (ii).

\subsection{$K3$ geometry, Mathieu moonshine, and the
$\operatorname{Sym}^{2}(L_{2})$ lock}\label{sec:k3}

\paragraph{Anchor (\tierB{}).} The elliptic genus of $K3$ decomposes into
characters whose multiplicities are dimensions of representations of the
Mathieu group $M_{24}$ (Eguchi--Ooguri--Tachikawa~\cite{EOT2011}); the
existence of the underlying $M_{24}$-module was proved by
Gannon~\cite{Gannon2016}. The mock-modular pieces are completed by shadows
following Zwegers~\cite{Zwegers2002}, and homological blocks connect such
mock objects to quantum invariants~\cite{GPPV2020}. For one-parameter
families of high Picard rank $K3$ surfaces (e.g.\ with Shioda--Inose
structure), the third-order Picard--Fuchs operator is the symmetric square
of a second-order operator~\cite{Doran2000}.

\paragraph{Conjecture (\tierC{}).}
\begin{conjecture}\label{conj:lock}
The macroscopic transport operator of the dual-scale system is conjugate to
$L_{3}=\operatorname{Sym}^{2}(L_{2})$, where $L_{2}$ is the Picard--Fuchs
operator of the fiber family; the non-holomorphic shadow completes the
T-dual layer to a modular object.
\end{conjecture}

\paragraph{Program note (\tierC{}, certificate pending).} Within Stream~1's
candidate selection over the sporadic sequences ($s_{7}$, $s_{10}$,
$S_{22}$), the $S_{12}$ sequence is \emph{provisionally} reclassified from
$K3$ candidate to elliptic-curve background. Under Rule~R5 this
classification remains Tier~C until its moduli map carries an
exact-arithmetic certificate of the same kind already obtained for the
(K-$s_{7}$, E-$s_{7/2}$) pairing.

\paragraph{Falsification.} A certified \textsc{fail} of the $S_{12}$ moduli
map against the elliptic-curve background, or a certified \textsc{pass}
against a $K3$ family, decides the reclassification either way.

\section{Reading the Lost Notebook: Motivation, Not Evidence}
\label{sec:lostnotebook}

The 1988 Narosa facsimile~\cite{Narosa1988} and the Andrews--Berndt volumes
\cite{AndrewsBerndt} organize a body of material with a genuine throughline
into the present program:

\begin{itemize}
\item The mock theta functions of the manuscript were completed to harmonic
Maass forms by Zwegers~\cite{Zwegers2002} and later connected to the $K3$
elliptic genus~\cite{EOT2011,Gannon2016}. That historical arc is real
(\tierB{}) and motivates Section~\ref{sec:k3}; it does not show that
Ramanujan was ``calculating string compactification dimensions.''
\item The continued-fraction material supplies the discrete-limit phenomena
(\tierB{}) that motivate Conjecture~\ref{conj:dim0}(ii).
\item The asymptotic material sits in the lineage of the circle method and
motivates the saddle-point entries of the dictionary.
\item The absence of connective proofs in the manuscript is a historical
observation about Ramanujan's working style. We adopt it as a \emph{design
metaphor} for the \textsc{Rama} engine. It is not evidence that his
cognition implemented the functional $E=\alpha C+\beta I+\gamma D$, and the
previous draft's claim to that effect is withdrawn.
\end{itemize}

\section{Formalization Status}\label{sec:lean}

\paragraph{Baseline.} The program's Lean~4 artifact
(\texttt{DualScale.lean}, rebuilt after the discovery of a vacuously-true
axiom derivation in an earlier file) satisfies Rules R1--R3: zero
\texttt{axiom} declarations and thirteen audit certificates. That baseline
is the standard against which this paper's targets are measured.

\paragraph{Withdrawal.} The five listings of the previous draft
(\texttt{HydrodynamicLimit}, \texttt{SpectralGap}, \texttt{CFMConstraints},
\texttt{PhaseTransitions}, \texttt{FiberLock}) all had the form
\texttt{theorem \ldots : True := by trivial}. Under Rule~R3 they certify
nothing and are withdrawn in full.

\paragraph{Honest targets.} The correct formal posture for an open
conjecture is an explicit, contentful statement marked \texttt{sorry}. The
following is \emph{schematic} (it does not typecheck as shown; the
supporting definitions are Stream~1 work items) and certifies nothing --- it
declares what would have to be closed for Conjecture~\ref{conj:U} to reach
Tier~A.

\begin{lstlisting}[caption={Formal target (OPEN). \texttt{sorry} = unproven; this listing is a specification, not a verification.}]
namespace DualScale.NS

/-- Truncation scale; exact rational, per audit rule R5. -/
def alphaPrime : Rat := 1 / 100

/-- Enstrophy density of the Galerkin truncation at |k| <= alphaPrime^(-1/2).
    (Definition to be supplied; specification only.) -/
-- def enstrophyDensity (u : TruncatedFlow alphaPrime) (t : Real) : Real := ...

/-- Hypothesis U, stated with content (uniformity in alphaPrime).
    OPEN TARGET: the `sorry` below is the honest marker of that fact. -/
theorem hypothesisU_uniform_bound :
    exists C : Real, 0 < C /\
      forall a : Rat, 0 < a -> a <= alphaPrime ->
        forall t : Real, a * enstrophyDensity_spec a t <= C := by
  sorry

end DualScale.NS
\end{lstlisting}

\noindent Audit note: any future claim that a target above has been closed
must be accompanied by a clean \texttt{\#print axioms} report and the
absence of \texttt{sorry} in the dependency cone.

\section{Milestones}\label{sec:milestones}

Each milestone has a decidable outcome.

\begin{description}
\item[M1.] Formalize the \emph{statements} of
Conjectures~\ref{conj:U}--\ref{conj:lock} so that all supporting definitions
compile; \texttt{sorry} only at the theorem level.
\item[M2.] Exact-arithmetic spectral certificates (Rule~R5) for the triad
graphs of Section~\ref{sec:gap} at $p\in\{2,3,5\}$: \textsc{pass} inside the
Alon--Boppana window or \textsc{fail} outside it.
\item[M3.] Leading-order derivation, or refutation, of the elliptic-angle
identity of Conjecture~\ref{conj:cfm}.
\item[M4.] Moduli-map certificate deciding the $S_{12}$ reclassification of
Section~\ref{sec:k3}.
\item[M5.] Conditional Tier~A theorem: Hypothesis~U $\Rightarrow$
uniform-in-$\alpha'$ smoothness on $[0,T]$ for the truncated family, with
the compactness step stated explicitly as a separate open lemma.
\end{description}

\section{Errata to Version 1}\label{sec:errata}

\begin{enumerate}
\item The author byline ``Fields Medal \& Nobel Prize in Physics Laureate
Framework'' is removed. No such distinction exists or is claimed.
\item All five Lean listings are withdrawn as vacuous
(\texttt{True := by trivial}); no machine verification of any mathematical
content occurred.
\item The phrases ``Hypothesis U is proven algebraically,'' ``singularities
are geometrically forbidden,'' and ``zero-axiom, computationally verified
bridge'' are retracted. Hypothesis~U is Conjecture~\ref{conj:U}; the
singular-set claim is Conjecture~\ref{conj:dim0} and is stronger than the
known CKN bound~\cite{CKN1982}.
\item The claim that the Lost Notebook ``proves that Ramanujan operated
precisely as a biological \textsc{Rama} engine'' is retracted; see
Section~\ref{sec:lostnotebook}.
\item The $S_{12}$ reclassification is downgraded from ``formally rejected''
to provisional, pending a Rule~R5 certificate.
\end{enumerate}

\section{Conclusion}

The worth of this program will be decided by which of
Conjectures~\ref{conj:U}--\ref{conj:lock} survive contact with certificates,
not by the reach of its vocabulary. Nothing here is currently proven. What
the program contributes today is a precise, falsifiable dictionary between
regularized hydrodynamics and the arithmetic of mock modular forms, and an
audit discipline --- no axioms, no vacuous goals, exact-arithmetic
certificates --- strict enough that when a Tier~A entry does appear, it will
mean something.

\begin{thebibliography}{99}\small

\bibitem{AJO2001}
G.~E. Andrews, J.~Jim\'enez-Urroz, K.~Ono,
\emph{$q$-series identities and values of certain $L$-functions},
Duke Math.\ J.\ \textbf{108} (2001), 395--419.

\bibitem{AndrewsBerndt}
G.~E. Andrews, B.~C. Berndt,
\emph{Ramanujan's Lost Notebook, Parts I--V},
Springer, New York, 2005--2018.

\bibitem{BKM1984}
J.~T. Beale, T.~Kato, A.~Majda,
\emph{Remarks on the breakdown of smooth solutions for the 3-D Euler
equations}, Comm.\ Math.\ Phys.\ \textbf{94} (1984), 61--66.

\bibitem{BowmanMcLaughlin2004}
D.~Bowman, J.~Mc\,Laughlin,
\emph{On the divergence of the Rogers--Ramanujan continued fraction on the
unit circle}, Trans.\ Amer.\ Math.\ Soc.\ \textbf{356} (2004).

\bibitem{CKN1982}
L.~Caffarelli, R.~Kohn, L.~Nirenberg,
\emph{Partial regularity of suitable weak solutions of the Navier--Stokes
equations}, Comm.\ Pure Appl.\ Math.\ \textbf{35} (1982), 771--831.

\bibitem{CF1993}
P.~Constantin, C.~Fefferman,
\emph{Direction of vorticity and the problem of global regularity for the
Navier--Stokes equations}, Indiana Univ.\ Math.\ J.\ \textbf{42} (1993),
775--789.

\bibitem{CFM1996}
P.~Constantin, C.~Fefferman, A.~Majda,
\emph{Geometric constraints on potentially singular solutions for the 3-D
Euler equations}, Comm.\ Partial Differential Equations \textbf{21} (1996),
559--571.

\bibitem{Deligne1974}
P.~Deligne,
\emph{La conjecture de Weil. I},
Publ.\ Math.\ IH\'ES \textbf{43} (1974), 273--307.

\bibitem{Doran2000}
C.~F. Doran,
\emph{Picard--Fuchs uniformization and modularity of the mirror map},
Comm.\ Math.\ Phys.\ \textbf{212} (2000), 625--647.

\bibitem{EOT2011}
T.~Eguchi, H.~Ooguri, Y.~Tachikawa,
\emph{Notes on the K3 surface and the Mathieu group $M_{24}$},
Exp.\ Math.\ \textbf{20} (2011), 91--96.

\bibitem{Gannon2016}
T.~Gannon,
\emph{Much ado about Mathieu},
Adv.\ Math.\ \textbf{301} (2016), 322--358.

\bibitem{GPPV2020}
S.~Gukov, D.~Pei, P.~Putrov, C.~Vafa,
\emph{BPS spectra and 3-manifold invariants},
J.\ Knot Theory Ramifications \textbf{29} (2020), 2040003.

\bibitem{LPS1988}
A.~Lubotzky, R.~Phillips, P.~Sarnak,
\emph{Ramanujan graphs},
Combinatorica \textbf{8} (1988), 261--277.

\bibitem{Narosa1988}
S.~Ramanujan,
\emph{The Lost Notebook and Other Unpublished Papers},
Narosa, New Delhi, 1988.

\bibitem{Zwegers2002}
S.~Zwegers,
\emph{Mock Theta Functions},
Ph.D.\ thesis, Universiteit Utrecht, 2002.

\end{thebibliography}

\end{document}